\documentclass[11pt]{article}

\usepackage{amsmath,amssymb}
\usepackage{xurl}
\usepackage[hidelinks]{hyperref}
\setlength{\emergencystretch}{2em}

% Shared commands for notes in docs/paper-gaps/.
% Notation follows the LDT paper (references/ldt-paper/).

\newcommand{\Lean}{\textsc{Lean}}
\newcommand{\leanid}[1]{\textsf{#1}}
\newcommand{\ghissue}[1]{\href{https://github.com/LionSR/MIPStarRE/issues/#1}{GitHub issue \##1}}
\newcommand{\ghpr}[1]{\href{https://github.com/LionSR/MIPStarRE/pull/#1}{GitHub PR \##1}}

% --- Paper-compatible notation ---
\newcommand{\F}{\mathbb{F}}
\newcommand{\N}{\mathbb{N}}
\newcommand{\C}{\mathbb{C}}
\newcommand{\tr}{\operatorname{tr}}
\newcommand{\ot}{\otimes}
\newcommand{\E}{\mathop{\bf E\/}}
\newcommand{\bx}{\mathbf{x}}
\newcommand{\by}{\mathbf{y}}
\newcommand{\bu}{\mathbf{u}}
\newcommand{\bv}{\mathbf{v}}

% Approximation relation (paper uses \approx_\eps and \simeq_\zeta)
\newcommand{\approxeps}{\approx_{\varepsilon}}
\newcommand{\simgeq}{\simeq_{\zeta}}

% Polynomial spaces (paper uses \polyfunc, \polysub, \polymeas)
\newcommand{\polyfunc}[3]{\operatorname{Poly}(#1,#2,#3)}
\newcommand{\polysub}[3]{\operatorname{PolySub}(#1,#2,#3)}
\newcommand{\polymeas}[3]{\operatorname{PolyMeas}(#1,#2,#3)}


\title{Proof-Gap Protocol for Source-Faithful Formalization}
\author{}
\date{2026-05-11}

\begin{document}
\maketitle

\section{Purpose}

This note fixes terminology for a recurring failure mode in the formalization
of the low individual degree test paper, \path{arXiv:2009.12982}.  It
complements the note-writing policy in \path{docs/paper-gaps/policy.tex} and
the source-statement audit tracked in \ghissue{1458}.
A proof may be blocked at
an intermediate mathematical assertion.  If that assertion is added as a new
hypothesis to a theorem labelled as a paper theorem, the resulting \Lean{}
theorem
may be correct as a conditional statement, but it is no longer the theorem in
the cited source.  The apparent proof has not closed the paper theorem; it has
moved a proof obligation into the theorem statement.

The protocol below distinguishes such proof debt from legitimate formal
boundary conditions and from genuine mathematical corrections.  It is meant for
contributors and reviewers who must decide whether a bridge, residual, repair,
input, or obligation structure is a proved local construction, a tracked proof
gap, or statement drift which must be repaired before a theorem may be
advertised as a source theorem.  The default rule is restrictive: such data
must not be added to a paper-facing theorem unless it is part of the cited
source statement or a documented faithful encoding of its ambient hypotheses.

\section{Basic objects}

\paragraph{Source theorem.}
A source theorem is a theorem, lemma, proposition, or corollary in the paper
source under \path{references/ldt-paper/}.  Its hypotheses and conclusion are
the mathematical ground truth, subject only to faithful formal encoding and to
documented paper-gap corrections.

\paragraph{Paper-facing declaration.}
A \Lean{} declaration is paper-facing when it is named or cited as the
formalization of a source theorem, for instance through a blueprint
\verb|\lean{}| link on a theorem environment carrying a paper label.  Such a
declaration must match the source theorem in its public statement.  In
particular, it may not add an unproved bridge, residual, repair, input,
obligation structure, or other proof-obligation hypothesis which the paper does
not assume.

\paragraph{Proof obligation.}
A proof obligation is an intermediate assertion needed by the proof of a source
theorem.  It may be a construction, an estimate, a transport lemma, or a
comparison between two formal representations.  If the paper proves or invokes
this assertion, then the formalization should state it as a separate lemma or
theorem and prove it from the upstream hypotheses.

\paragraph{Obligation discharger.}
An obligation discharger is a theorem or definition which constructs a proof
obligation from the paper hypotheses already available at that point in the
argument.  This is the preferred representation of missing work.  If the proof
is not yet complete, the discharger may temporarily contain a tracked
\verb|sorry|, but the paper-facing theorem should call it rather than assume
its output as an additional theorem hypothesis.

\paragraph{Conditional helper.}
A conditional helper is a separately named theorem whose statement explicitly
assumes a proof obligation.  Names such as
\leanid{mainFormal\_ofInternalObligations},
\leanid{orthonormalization\_ofInput}, or names whose suffix advertises the
assumed obligation are permitted only when the declaration is not presented as
the source theorem.  A conditional helper is a quarantine for downstream proof
content; it does not close the source theorem.

\section{Bridge, residual, repair, input, and obligation structures}

The words \emph{bridge}, \emph{residual}, \emph{repair}, \emph{input}, and
\emph{obligation structure} name possible formal shapes.  They become dangerous
when their mathematical status is unclear or when their names make unproved
data look like acceptable hypotheses of a source theorem.

\paragraph{Bridge.}
A bridge should be a transport theorem between two already defined
mathematical forms: for example, a tensor-placement identity, a conversion
between two equivalent consistency relations, or a statement translating a
source object into a formal representation.  A bridge is harmless only after it
is proved from previously available assumptions.  It becomes proof debt when it
is assumed because the transport has not yet been proved.  It becomes theorem
drift when the assumed bridge is added to a paper-facing theorem.

\paragraph{Residual.}
A residual is a named remainder after part of an argument has been performed.
It is often a useful way to record the exact missing target after a completed
construction.  A residual may appear as the conclusion of a helper theorem or
as the input of a quarantined conditional helper.  It is not a permissible new
assumption on a paper theorem unless the paper itself assumes that residual.

\paragraph{Input.}
An input is a named bundle of data or hypotheses consumed by a construction or
conditional helper.  It is harmless only when each load-bearing field is either
part of the source hypotheses or is produced by an earlier theorem.  It becomes
proof debt when the bundle is merely assumed because the construction of its
fields is missing.  It becomes theorem drift when such a bundle is added to the
public statement of a source theorem.

\paragraph{Repair.}
A repair is a mathematical construction which replaces an approximate or
partial object by a better object, usually with a quantitative estimate.  In
the low individual degree test formalization, projectivization and
orthonormalization repairs are genuine mathematical operations.  A repair
should therefore appear as a construction or as a theorem discharging the
corresponding obligation.  If a theorem assumes the output of the repair instead
of constructing it, that theorem is conditional on the repair.

\paragraph{Obligation structure.}
An obligation structure is a record whose fields are the data and estimates
needed by a later construction.  This is acceptable only as an internal target
or as a hypothesis of a conditional helper.  Such a structure becomes statement
drift if it is added to a paper-facing theorem instead of being constructed
inside the proof.  The review question is whether every load-bearing field is
proved from the source hypotheses or is explicitly tracked as a missing
obligation.

\section{The main failure mode}

Suppose the paper proves
\[
  H \Longrightarrow C.
\]
During formalization, the proof may need an intermediate assertion \(B\).  The
paper theorem should be represented as
\[
  H \Longrightarrow B \Longrightarrow C
  \quad\text{inside the proof,}
\]
where \(B\) is obtained by a lemma or theorem from \(H\).  The following
public statement is a different theorem:
\[
  H \land B \Longrightarrow C.
\]
It is not unsound in \Lean{}, but it is not the paper theorem.  It says that the
conclusion follows if the missing bridge data is supplied.  Such a theorem is a
conditional helper, even if its conclusion looks exactly like the paper's
conclusion.

This is the situation colloquially described as ``adding a bridge
hypothesis''.  More precisely, an unproved intermediate step has been promoted
from the proof body to the theorem's assumptions.  The correction is not to hide
the hypothesis under a different structure name.  The correction is to keep a
paper-facing theorem with the source statement and move the missing step to a
named internal proof obligation.

\section{Faithful boundary hypotheses}

Not every extra formal hypothesis is a theorem drift.  Some hypotheses are
faithful boundary conditions introduced by formal encoding.  Examples include
nonnegativity needed for a real power, positivity needed to divide by a
quantity, finiteness and nonemptiness of an index type, a field-model instance,
or decidability instances required by \Lean{}.  Such hypotheses are acceptable
when they express part of the mathematical ambient setting or make explicit a
condition used implicitly in the paper.

The distinction is load-bearing.  A faithful boundary hypothesis describes the
domain in which the paper's objects live.  A proof-debt hypothesis supplies a
mathematical step that the proof should derive.  If removing the hypothesis
leaves an unproved construction or estimate from the paper proof, then the
hypothesis is proof debt, not a boundary condition.

\section{Definition drift and subtree repair}

Definition drift is especially dangerous because it propagates.  If a
definition is changed so that it no longer denotes the object in the paper,
then every downstream theorem using that definition may continue to compile
while proving statements about the wrong object.  Downstream theorem statements
can then be repaired piecemeal, but the development remains unstable because
the wrong definition is still upstream.

When a definition mismatch is found, the repair should proceed by a subtree
protocol.

\begin{enumerate}
  \item Identify the source object and the formal object precisely.
  \item Freeze a source-faithful definition or state an adapter theorem
  relating the existing definition to the source object.
  \item Recheck the downstream theorem chain against that definition or adapter.
  \item Re-formalize the affected subtree from the corrected interface, rather
  than adding local hypotheses to each downstream theorem.
  \item If a downstream proof cannot be completed, record the missing obligation
  at the lowest point where the paper proof needs it.
\end{enumerate}

This is slower than editing a final theorem signature, but it prevents a local
definition error from becoming a global theorem-statement error.

\section{Allowed shapes}

\begin{center}
\begin{tabular}{p{0.24\linewidth}p{0.31\linewidth}p{0.32\linewidth}}
\textbf{Shape} & \textbf{Allowed use} & \textbf{Forbidden use} \\
\hline
Paper-facing theorem & Matches the source statement up to faithful encoding &
Carries non-paper bridge, residual, repair, input, obligation-structure, or
proof-obligation hypotheses \\
Conditional helper & Assumes a named obligation and has a conditional name &
Is linked as the source theorem or marked as proof-complete for that theorem \\
Obligation discharger theorem & Constructs the obligation from paper hypotheses &
Merely restates the obligation as an assumption \\
Input bundle & Collects produced data for a local construction & Is assumed by a
source theorem instead of being produced \\
Obligation structure & Bundles constructed data and estimates & Hides unproved
fields that are then assumed by a source theorem \\
Paper-gap note & Documents a real mathematical or formal-statement divergence &
Justifies silent statement drift without a repair plan
\end{tabular}
\end{center}

\section{Protocol for a blocked proof}

When a source-labelled proof is blocked, use the following order.

\begin{enumerate}
  \item Read the cited source statement and the proof passage in
  \path{references/ldt-paper/}.
  \item Write the paper theorem form, the current \Lean{} theorem form, and the
  first missing intermediate assertion.
  \item Decide whether the missing assertion is a faithful boundary condition,
  a representation adapter, an ordinary lemma from the paper proof, or a real
  paper-gap discrepancy.
  \item If it is an ordinary lemma or adapter, create a named theorem which
  discharges that obligation.  Do not add it to the source theorem as a new
  hypothesis.
  \item Keep the source-labelled theorem itself in the source form, even if its
  proof must contain a tracked \verb|sorry|.  This is the default repair when
  the missing assertion has not yet been proved.
  \item Create a conditional helper only to quarantine substantial proof content
  that already exists and would otherwise be lost.  Do not create such a helper
  merely to avoid a \verb|sorry| or to keep making progress.  If a conditional
  helper remains, give it a conditional name, keep it separate from the
  paper-facing theorem, and record the obligation discharger that must eliminate
  it.
  \item If the source and formal statements genuinely differ, write or update a
  note under \path{docs/paper-gaps/} before marking the statement as aligned.
  \item Update the blueprint so that paper labels point only to paper-facing
  declarations.  Conditional helpers belong in remarks or implementation
  notes, not in the source theorem block.
\end{enumerate}

\section{Review questions}

A reviewer should first compare every paper-labelled declaration with the cited
source statement, checking both hypotheses and conclusion up to faithful formal
encoding.  Bridge, residual, repair, input, obligation-structure, and
proof-obligation assumptions should be absent from paper-facing theorem
signatures unless they occur in the paper.  If a conditional helper remains,
its name should be explicitly conditional and it should not support a source
theorem \verb|\leanok| claim.
Each load-bearing obligation-structure or input field should have a theorem
discharging it, a tracked issue, or an explicit \verb|sorry| at the obligation
site.  When a definition has changed, the downstream subtree should be checked
against the corrected definition rather than patched locally.  Finally, if the
formal statement genuinely differs from the source statement, a paper-gap note
should explain the mathematical difference and give a plan to remove or justify
it.

\section{Conclusion}

The project should prefer visible proof obligations over invisible statement
drift.  A tracked \verb|sorry| in the source-shaped theorem, or in the lowest
named internal obligation needed by its proof, is less misleading than a source
theorem whose assumptions have been strengthened.  Conditional helpers are not a
normal substitute for unfinished proofs; they are temporary quarantines for
substantial proof content, and they must remain conditional until the missing
obligation is discharged.  The paper-facing declaration is reserved for the
theorem proved by the paper, or for a documented corrected version whose
mathematical difference is explicit.

\end{document}
